%% =================================================================
%% Wei, Mei, Zhao, Xiao, Sun, Zhang, Guo.
%% "Nondestructively identifying the mechanical behavior of soft
%%  tissues using surface deformation with an explicit inverse
%%  approach."
%% Appl. Math. Modelling 134 (2024) 126--147.
%%
%% Reconstruction of page 20 (printed p. 145) as study notes.
%% Generated by: reproduce-basic-paper-tex → reproduce-multiple-pages-basic-tex [17, 22]
%%
%% Structure: Fig. 26 (Sample 2 two-inclusion, no noise) + Table 6
%% (regularization factors for Figs 23-26) + prose transition
%% (future-research paragraph + overall summary) + Conclusion heading
%% with closing paragraph. No equations.
%% =================================================================

\documentclass[11pt]{article}

\usepackage[margin=1in]{geometry}
\usepackage{amsmath, amssymb}
\usepackage{mathrsfs}
\usepackage{bm}
\usepackage{xcolor}
\usepackage{fancyhdr}
\usepackage{booktabs}
\usepackage{multirow}
\usepackage{array}

\pagestyle{fancy}
\fancyhf{}
\lhead{\small Y.~Mei et al.}
\rhead{\small Applied Mathematical Modelling 134 (2024) 126--147}
\cfoot{\small 145 $\to$ \arabic{page}}
\renewcommand{\headrulewidth}{0.4pt}

\setcounter{page}{1}

\begin{document}

% ---------- Figure 26 placeholder ----------
\noindent
\begin{center}
\fbox{\parbox{0.92\textwidth}{\centering\vspace{1em}
\textbf{[Figure 26 --- placeholder]}\\[0.8em]
Six-panel comparison of \textbf{Explicit Inverse Method} (top
row) vs.\ \textbf{Implicit Inverse Method} (bottom row) on
\textbf{Case 3 Sample 2} (two embedded inclusions side-by-side,
target $\approx 5$\,MPa each), with \textbf{no noise}, using 3,
6, and 9 displacement datasets.\\[0.4em]
\textbf{Explicit (top row)} --- both inclusions are recovered as
distinct red blobs on a uniform blue background. Colorbar
ranges:
(a) 3 fields, SM $\in [1.00, 7.06]$\,MPa --- \textbf{41\,\%
overshoot};
(b) 6 fields, SM $\in [1.00, 6.06]$\,MPa --- 21\,\% overshoot;
(c) 9 fields, SM $\in [1.00, 5.97]$\,MPa --- 19\,\% overshoot.
Max drifts downward monotonically as more fields are added,
approaching but not reaching the 5\,MPa nominal.\\[0.4em]
\textbf{Implicit (bottom row)} --- both inclusions are
recovered as green/yellow blobs (distinctly underbright relative
to the explicit method's red). Colorbar ranges:
(d) 3 fields, SM $\in [1.00, 2.61]$\,MPa --- \textbf{48\,\%
undershoot};
(e) 6 fields, SM $\in [1.00, 3.00]$\,MPa --- 40\,\% undershoot;
(f) 9 fields, SM $\in [1.00, 3.82]$\,MPa --- 24\,\% undershoot.
Max drifts upward monotonically, approaching but nowhere close
to the 5\,MPa nominal.\\[0.3em]
\textbf{Takeaway}: on the two-inclusion case (even without
noise), the explicit method overshoots by 19--41\,\% and the
implicit method undershoots by 24--48\,\%. The paper's prose
reports both deviations (see page~18: ``overestimated by
approximately 20\,\%'' for explicit, ``underestimation ...\ by
at least 20\,\%'' for implicit). In this case the prose framing
and the visible data agree --- unlike the Fig.~25 case where
the ``2\,\% average error'' claim hid a 7\,\% inclusion-specific
undershoot.
\vspace{1em}}}
\end{center}
\vspace{0.3em}
\begin{center}
\textbf{Fig. 26.} The reconstructed results with varying 3, 6
and 9 surface displacement datasets corresponding to Sample 2
(Without noise, Unit: MPa).
\end{center}

\vspace{1em}

% ---------- Table 6 ----------
\noindent\textbf{Table 6}\\[0.2em]
{\small The regularization factors of Figs.~23--26.}
\vspace{0.3em}

\begin{center}
\footnotesize
\begin{tabular}{@{}ll ccc ccc@{}}
\toprule
\multicolumn{2}{l}{Inverse method}
  & \multicolumn{3}{c}{Explicit inverse method}
  & \multicolumn{3}{c}{Implicit inverse method} \\
\cmidrule(lr){3-5}\cmidrule(lr){6-8}
\multicolumn{2}{l}{Number of surface displacement datasets}
  & 3 & 6 & 9 & 3 & 6 & 9 \\
\midrule
\multirow{3}{*}{Sample 1}
  & no noise    & $1\times 10^{-15}$ & $1\times 10^{-16}$ & $1\times 10^{-17}$ & $1\times 10^{-14}$ & $1\times 10^{-14}$ & $1\times 10^{-14}$ \\
  & 1\,\% noise & $1\times 10^{-13}$ & $1\times 10^{-13}$ & $2\times 10^{-13}$ & $5\times 10^{-13}$ & $3\times 10^{-13}$ & $2\times 10^{-12}$ \\
  & 5\,\% noise & $1\times 10^{-12}$ & $1\times 10^{-12}$ & $1\times 10^{-12}$ & $1\times 10^{-11}$ & $1\times 10^{-11}$ & $1\times 10^{-12}$ \\
\midrule
Sample 2 & no noise & 0.00 & 0.00 & 0.00 & $1\times 10^{-13}$ & $1\times 10^{-13}$ & $1\times 10^{-14}$ \\
\bottomrule
\end{tabular}
\end{center}

\vspace{1em}

% ---------- Prose transition ----------
\noindent
future research, we plan to explore the integration of optical
measurement techniques with the proposed explicit inverse
approach to conduct practical experiments and further test this
novel structural health monitoring approach. By combining
optical measurement data with our method, we aim to validate
its effectiveness in real-world scenarios and enhance its
applicability for monitoring the variation of biomechanical
behavior of soft tissues. This would provide valuable insights
into the method's performance under realistic conditions and
potentially open up new avenues for its practical implementation
in biomedical applications.

Overall, the explicit inverse approach presented in the paper
offers a more efficient and cost-effective method for
characterizing the elastic property distribution of biological
structures with surface deformation. The reduction in
optimization variables and the reliance on easily measurable
surface deformation data make it a valuable contribution to the
field.

\vspace{0.5em}

\noindent\textbf{Conclusion}

\vspace{0.3em}

\noindent
In this paper, we propose an explicit inverse approach to map
the nonhomogeneous elastic property distribution of three
typical tissue structure by surface deformation nondestructively.
Several numerical examples are employed to assess the
feasibility and performance of the proposed method. The results
of the numerical study demonstrate the robust capability of the
explicit inverse method in accurately reconstructing the
nonhomogeneous shear modulus distribution of soft tissues,
surpassing the performance of the implicit inverse method with
limited surface deformation datasets. Specifically for the
layered ring structure, we observe that the average relative
error of the reconstruction can exceed 40\,\% when using the
implicit inverse method, while the explicit inverse method
achieves a remarkable average relative error of only 5\,\%.
The outcomes of our research underscore the efficacy of the
proposed explicit inverse approach, offering a more efficient
and cost-effective means to characterize the elastic property
distribution of biological structures utilizing a reduced
number of optimization variables.

% [page ends here — continues on page 21]

\end{document}
